% The preamble of the two roots of the theory paper (main.tex, supplement.tex), on the model of
% the isotropic paper's (spatial-hemigroup-affine, paper-iso/preamble.tex). A root that defines
% \pyrexternal as the name of the other document reads that document's labels through xr-hyper,
% which is loaded before hyperref; the option nocite keeps the other document's bibliography
% numbers out, so that each has its own list of references.
\usepackage[a4paper,margin=1in]{geometry}
\setlength{\emergencystretch}{1.5em}
\usepackage{amsmath,amssymb,amsthm}
\usepackage{booktabs}
\usepackage{array}
\usepackage{longtable}
\newcolumntype{R}[1]{>{\raggedright\arraybackslash}p{#1}}
\newcolumntype{S}[1]{>{\footnotesize\raggedright\arraybackslash}p{#1}}
\setlength{\LTcapwidth}{\textwidth}
\usepackage{graphicx}
\usepackage{tikz}
% The figures are drawn by scripts/make_figures_theory.py into figures/ (theory-*.png).
\graphicspath{{../figures/}}
% natbib in numeric mode, for \citenum: the blueprint writes a citation with its locator as
% [\sscite{key}, Thm.~x, p.~y], so \sscite prints the bare number (macros.tex).
\usepackage[numbers]{natbib}
\ifdefined\pyrexternal\usepackage{xr-hyper}\fi
\usepackage[hidelinks]{hyperref}
\hypersetup{
  pdftitle={Spatial Scale Space from Hemigroup Axioms: The Theory of Image Pyramids},
  pdfauthor={Daniel Fagerstr\"om},
}

% The macros of the theory paper and its supplement. The blueprint's (blueprint/src/macros.tex)
% are kept by name, so that a statement shared with the blueprint compiles here unchanged and
% `linkage check` compares the same source text; three of them print differently here:
%   \prog{<module> v0.1}{<number>}  a citation of the released module, by its printed number;
%   \ledger{A<n>}                   the cited fact's own statement in section 2.2, by number;
%   \sscite{<key>}                  \cite.
% The leanblueprint metadata (\lean, \uses, \leanok, \notready) is not transcribed; the commands
% are defined empty so that a statement copied with them would still build.

\newcommand{\RR}{\mathbb{R}}
\newcommand{\CC}{\mathbb{C}}
\newcommand{\NN}{\mathbb{N}}
\newcommand{\ZZ}{\mathbb{Z}}
\newcommand{\Cin}{\operatorname{Cin}}   % the entire cosine integral, as in the cone module
\newcommand{\Levy}{L\'evy}              % the name, accented through one macro
\newcommand{\Matern}{Mat\'ern}          % the same device for the Matern family
\newcommand{\Polya}{P\'olya}            % and for Polya frequency functions
\newcommand{\GL}{\mathrm{GL}}
\newcommand{\Orth}{\mathrm{O}}          % the orthogonal group O(d)
\newcommand{\tp}{\mathsf{T}}            % transpose, written A^{\tp}
\newcommand{\tr}{\operatorname{tr}}
\newcommand{\diag}{\operatorname{diag}}
\newcommand{\PD}{\mathrm{PD}}           % positive definite d x d matrices
\newcommand{\Id}{\mathrm{Id}}
\newcommand{\Geom}{\mathrm{Geom}}       % the geometric law

\providecommand{\lean}[1]{}
\providecommand{\leanok}{}
\providecommand{\uses}[1]{}
\providecommand{\notready}{}
\providecommand{\proves}[1]{}
\providecommand{\notes}[1]{}
\providecommand{\statedin}[1]{}
\providecommand{\restates}[1]{}

% The blueprint writes every citation as [\sscite{key}] or [\sscite{key}, <locator>], so \sscite
% prints the bare number and the brackets are the text's.
\newcommand{\sscite}[1]{\citenum{#1}}

% The released modules of the programme, by the first argument of \prog as the blueprint writes
% it; the citekeys are the library's (hub constellation.json; paper/references.bib).
\makeatletter
\@namedef{progkey@V-line v0.1}{fagerstrom2026spatial}
\@namedef{progkey@V-cone v0.1}{fagerstrom2026cone}
\@namedef{progkey@V-selection v0.1}{fagerstrom2026selection}
\@namedef{progkey@iso v0.1}{fagerstrom2026isotropic}
\@namedef{progkey@affine v0.1}{fagerstrom2026affine}
\newcommand{\prog@key}[1]{\@ifundefined{progkey@#1}{\PackageError{macros}{No citation for the
  programme label #1}{}}{}\edef\prog@thekey{\csname progkey@#1\endcsname}}
\newcommand{\prog@cite}[2]{\cite[#2]{#1}}
\DeclareRobustCommand{\prog}[2]{\prog@key{#1}\expandafter\prog@cite\expandafter{\prog@thekey}{#2}}
\DeclareRobustCommand{\progpaper}[1]{\prog@key{#1}\expandafter\cite\expandafter{\prog@thekey}}
\makeatother

% A ledger entry is printed as its statement in section 2.2 (the node thm:cited-A<n>). The \ref is
% reached through \csname so that `linkage paper`, which reads \ref{...} arguments as labels, does
% not take the macro parameter for one.
\newcommand{\ledger}[1]{Theorem~\csname ref\endcsname{thm:cited-#1}}

% The other document is imported after macros.tex: its .aux carries titles written in the
% paper's macros (\Levy), which xr expands as it reads them.
\ifdefined\pyrexternal\externaldocument[][nocite]{\pyrexternal}\fi
\numberwithin{equation}{section}

% The blueprint numbers its four statement environments on one counter; so does the paper.
\theoremstyle{plain}
\newtheorem{theorem}{Theorem}[section]
\newtheorem{proposition}[theorem]{Proposition}
\newtheorem{lemma}[theorem]{Lemma}
\newtheorem{corollary}[theorem]{Corollary}
\theoremstyle{definition}
\newtheorem{definition}[theorem]{Definition}
\newtheorem{example}[theorem]{Example}
\theoremstyle{remark}
\newtheorem{remark}[theorem]{Remark}
